\documentclass[10pt,twocolumn,letterpaper]{article}
\usepackage[pagenumbers]{cvpr}
% Tectonic/XeTeX uses TU encoding; map CVPR Times to installed Times New Roman.
% Standard pdfLaTeX keeps cvpr.sty's normal Times configuration.
\ifdefined\XeTeXversion
\DeclareFontFamily{TU}{ptm}{}
\DeclareFontShape{TU}{ptm}{m}{n}{<-> "TimesNewRomanPSMT"}{}
\DeclareFontShape{TU}{ptm}{b}{n}{<-> "TimesNewRomanPS-BoldMT"}{}
\DeclareFontShape{TU}{ptm}{m}{it}{<-> "TimesNewRomanPS-ItalicMT"}{}
\DeclareFontShape{TU}{ptm}{b}{it}{<-> "TimesNewRomanPS-BoldItalicMT"}{}
\DeclareFontFamily{TU}{pcr}{}
\DeclareFontShape{TU}{pcr}{m}{n}{<-> "CourierNewPSMT"}{}
\fi
\usepackage{amsmath,amssymb,booktabs}
\usepackage[breaklinks,colorlinks,citecolor=blue,urlcolor=blue,linkcolor=blue]{hyperref}
\captionsetup{hypcap=false}
\urlstyle{same}
\title{Clinical Trial Stopping Reasons and Outcome Prediction\\
\large CS5785 / ORIE5750 / ECE5414: Fall 2026 Midterm Project}
\hypersetup{pdfauthor={Pei-Chen Chung (pc786) and Yu-Jye Hu (yh2444)}}
\author{Pei-Chen Chung (NetID: pc786)\qquad Yu-Jye Hu (NetID: yh2444)\\Cornell University\\
\small Both authors contributed equally to data preparation, analysis, modeling, evaluation, and writing.}
\begin{document}
\maketitle
\begin{abstract}
We study scarce-label clinical-trial stopping reasons and phase-specific outcomes.
K-Means and GMM recover masked reason labels, but gains depend on the classifier and decision policy.
Calibrated RBF kernel ridge raises reported public sample F1 from 0.85405 to 0.88568;
a phase-aware text blend lowers public log loss from 0.60149 to 0.59011.
Grouped and temporal diagnostics, feature ablations, and reason-specific errors qualify these gains.
\end{abstract}
\section{Methods}
\subsection{Data and evaluation}
The supplied data define both tasks~\cite{handout}. Part 1 has 1,171 training trials:
157 observed five-label vectors and 1,014 masked as $-100$; validation/test have 78/312 trials.
Gold seeds oversample early trials and rare reasons. Validation spans 1999--2009;
test spans 2010--2016. IDs only identify/join records, never predict labels.

Sample F1 averages across trials; per-reason and macro F1 are also reported, with undefined values zero.
We distinguish validation selection from later label reuse. Final decisions threshold all reasons,
then use argmax if empty; ties follow Safety, Funding, Efficacy, Recruitment, Other.

\paragraph{Part 1 feature choice.}
Part 1 retains all 150 reason-embedding coordinates, summarizing stopping-reason text.
PCA/random rotation remove coordinate interpretation, so we compare feature blocks and scaling.
Final ridge uses embeddings alone; year defines temporal diagnostics.
Pathway and text additions are evaluated in Section~\ref{sec:p1features}.

\subsection{Label recovery and supervised controls}
Five observed-label logistic regressions (LR) form the supervised baseline, with $C=1$ and threshold 0.5.
StandardScaler fits available training embeddings. K-Means uses 10 clusters, k-means++, 20 starts,
Lloyd updates, 500 maximum iterations, and tolerance $10^{-4}$; the recorded fit takes 19 iterations.
GMM uses 10 diagonal components, five K-Means initializations, 500 maximum EM iterations,
tolerance $10^{-3}$, and covariance regularization $10^{-4}$; it converges in 23 iterations.
For cluster memberships/responsibilities $w_{ik}$ and labeled rows $L$,
\begin{equation}
 q_{kr}=\frac{\sum_{i\in L}w_{ik}y_{ir}+2\pi_r}{\sum_{i\in L}w_{ik}+2},\qquad
 \pi_r=|L|^{-1}\sum_{i\in L}y_{ir}.
\end{equation}
Masked labels are $\mathbb{I}[\sum_k w_{ik}q_{kr}\geq0.5]$; empty clusters use the prior.
Gold targets stay unchanged. Gold/pseudo weights are 1/0.1~\cite{lectures}.
LR uses liblinear, 500 iterations, tolerance $10^{-4}$, and no class weights~\cite{sklearn}.
We compare $C=0.1,1,10$, thresholds 0.35/0.5/0.65, and broad pathway indicators.
The historical choice is K-Means plus LR with $C=0.1$, threshold 0.35.

\subsection{Local models and final Part 1 fit}
The later study compares 24 configurations: column standardization, global scaling, or row L2 normalization~\cite{normalizer};
RBF SVM~\cite{svm} or kernel ridge; two bandwidths and two regularization settings.
The winning full-data recipe preserves Euclidean geometry through global scaling and uses
\begin{equation}
 \begin{aligned}
 K_{ij}&=\exp\!\left(-0.5\frac{\|x_i-x_j\|^2}{m}\right),\\
 A&=(K+0.1I)^{-1}Y,
 \end{aligned}
\end{equation}
where $m$ is the median nonzero squared distance among fitting gold rows.
Geometry uses eligible gold and unlabeled embeddings (235+1,014 in the final fit);
regression and calibration use gold targets only.
Five sigmoid calibrators use gold-only, three-fold grouped OOF scores, standardized
before LR calibration ($C=1$). Predictions use threshold 0.5 plus the declared fallback.
After selection, this supervised model fits all 235 visible labels. Its original validation labels
are then training data; no independent validation score is claimed for that final fit.

\clearpage
\twocolumn[{
\section{Results: Stopping Reasons}
\begin{center}\small
\begin{tabular}{lrrrrrrr}
\toprule
Model & Sample F1 & Macro F1 & Safety & Funding & Efficacy & Recruit. & Other \\
\midrule
Supervised LR & 0.8782 & 0.6990 & 0.8571 & 0.2857 & 0.8125 & 0.9796 & 0.5600 \\
Initial K-Means + LR & 0.8547 & 0.6690 & 0.8235 & 0.4000 & 0.7407 & 0.9524 & 0.4286 \\
Retained K-Means + LR & 0.8782 & 0.7648 & 0.8421 & 0.6667 & 0.8387 & 0.9434 & 0.5333 \\
Fixed ridge recipe (157 gold) & 0.8932 & 0.7829 & 0.9412 & 0.6667 & 0.8000 & 0.9804 & 0.5263 \\
\bottomrule
\end{tabular}

\captionof{table}{Historical validation on 78 trials, fitted using 157 original gold training rows.
All rows use thresholding followed by argmax if empty. These are pre-refit diagnostic comparisons;
the final public ridge model subsequently uses all 235 visible labels. Positive support is
9 Safety, 4 Funding, 14 Efficacy, 50 Recruitment, and 10 Other.}
\label{tab:p1val}
\end{center}
\begin{center}\small
\begin{tabular}{lrrr}
\toprule
Evaluation & Fixed K-Means + LR & Local selection procedure & Fixed ridge recipe\\
\midrule
Repeated near-group OOF, 235 labels & 0.8877 & 0.8983 & 0.9026\\
Historical validation, 78 trials & 0.8782 & 0.8803 & 0.8932\\
Late train, 20 trials & 0.8417 & 0.8833 & 0.9000\\
Late combined, 56 trials & 0.8393 & 0.8869 & 0.8512\\
\bottomrule
\end{tabular}
\captionof{table}{Sample F1 under explicitly different procedures. The middle column evaluates
nested selection among local models, consensus students, and blends. The right column was
measured later for the already-selected public-winning ridge recipe; it is not an independent
estimate obtained before choosing that recipe. Temporal fits use 134 early train labels or
176 early train/validation labels respectively. The same labeled and unlabeled held groups
are excluded in OOF; temporal diagnostics retain earlier matching groups.}
\label{tab:p1eval}
\end{center}
}]
\subsection{What changed the conclusion}
Baseline F1 0.8654 omitted four empty predictions; the submission fallback raises it to 0.8782,
tying retained K-Means LR. Validation alone establishes no semi-supervised F1 gain.
Exact validation sets are 61/78 versus 64/78; test sets differ on 96/312 trials.
Earlier 157-label grouped OOF favored K-Means (0.8875 versus 0.8399), under different
grouping and populations from Table~\ref{tab:p1eval}.

Under the original empty-set policy, K-Means/GMM with LR $C=1$ score 0.8205/0.6667.
GMM produces no Safety, Funding, or Efficacy pseudo positives; K-Means produces no Funding
positives and only 39 Other positives among 1,014 masked rows. This label-quality limitation
motivates local models without inherited cluster targets.

\subsection{Grouped selection and calibration}
The 235 visible labels form 166 near-duplicate groups: connected components at raw Euclidean
distance $\leq10^{-4}$ across gold and masked rows. Held groups leave all fitting pools.
Four outer folds repeat with seeds 42, 137, and 2026; three inner group folds select models
and decisions, with further grouped calibration inside each fit. The 705 OOF predictions
repeat 235 trials, rather than represent independent observations.

When the local model and K-Means/LR control have complementary inner errors, local and
gold-only linear teachers can
supply label-wise soft pseudo targets. All six member scores must be at least 0.8 or at most
0.1; otherwise that label remains unknown. Gold weights remain one, pseudo weights are 0.1
with separate class-support caps, and only one student round is fitted. No outer fold selected
this student. Three of 12 folds selected blends; nine selected local models alone.

\subsection{Observed improvement and remaining errors}
Full-data selection chose ridge; 11 outer folds selected SVMs. Its public F1 is 0.88568,
up 0.03163 from K-Means LR. This aggregate feedback reveals no individual hidden labels.
The procedure's paired group-bootstrap 95\% gain interval over the inner-retuned control,
$[-0.00690,0.03494]$, includes zero.

The procedure improves Safety OOF F1 from 0.8228 to 0.9103, while Other declines from 0.5865
to 0.5761. Nine distinct Recruitment+Other trials are predicted as Recruitment alone in all
three repeats. The main improvement therefore does not resolve this rare co-occurrence.

\subsection{Feature ablation}
\label{sec:p1features}
With K-Means pseudo labels, gold/pseudo weights, LR $C=1$, and threshold 0.5 fixed,
we append 27 level-0 Reactome pathway indicators (training frequency $\geq5$) and two
availability flags to the 150 embeddings. Validation sample/macro F1 falls from
0.8205/0.6332 to 0.8120/0.5680. This controlled feature comparison uses the original
empty-set-allowing policy, unlike Table~\ref{tab:p1val}.
A separate TF-IDF view of criteria, diseases, and drugs supplies conservative Other additions;
it changes no decisions at its OOF-selected threshold: validation F1 remains 0.8782 and
late-56 F1 0.8393. This tests a targeted correction, not all text/embedding combinations.

\clearpage
\twocolumn[{
\section{Results: Outcome Prediction}
\begin{center}\small
\begin{tabular}{lrr}
\toprule
Model / feature representation & Validation equal-phase log loss & Reported public log loss\\
\midrule
Phase-specific word TF-IDF LR baseline & 0.58377 & 0.60149\\
50/50 concatenated-word / separate-field blend & 0.58227 & 0.59962\\
25/25/50 word / field / character blend & 0.58090 & 0.59696\\
50/50 historical blend / phase-shared text model & \textbf{0.56839} & \textbf{0.59011}\\
\bottomrule
\end{tabular}
\captionof{table}{The final outcome model retains the historical text blend as one component.
Lower loss is better. Public values come from user-provided Kaggle screenshots; private results
are unknown. Historical delivery files remain separately reproducible.}
\label{tab:p2models}
\end{center}
\begin{center}\small
\begin{tabular}{llrrrr}
\toprule
Model & Phase & F1 & ROC-AUC & AP & Log loss \\
\midrule
Baseline & I & 0.7795 & 0.7820 & 0.8285 & 0.5966 \\
Baseline & II & 0.6236 & 0.7236 & 0.7088 & 0.6082 \\
Baseline & III & 0.8126 & 0.7663 & 0.8685 & 0.5464 \\
Current phase-shared blend & I & 0.7742 & 0.7961 & 0.8400 & 0.5738 \\
Current phase-shared blend & II & 0.6297 & 0.7387 & 0.7259 & 0.5983 \\
Current phase-shared blend & III & 0.8178 & 0.7817 & 0.8776 & 0.5330 \\
\bottomrule
\end{tabular}

\captionof{table}{Required validation metrics by phase for the original baseline and current
public-best outcome model. Success is the positive class; F1 uses 0.5, ROC-AUC and average
precision (AP) use probabilities. Validation has 149/501/363 phase-I/II/III rows.}
\label{tab:p2phase}
\end{center}
}]
\subsection{Models and feature choices}
Part 2 has 8,037 training and 1,013 validation trial-phase rows. All phase copies of a trial stay
together in model-selection splits. Test has 1,500 rows, 500 per phase, so equal-phase mean
log loss is the selection objective. Observed test start years are 2014--2021, later than the
training/validation population through 2013.

The baseline fits one TF-IDF LR per phase on concatenated eligibility criteria, disease, and
drug text~\cite{tfidf,sklearn}. Word unigrams/bigrams use minimum document frequency 3,
at most 20,000 features, lowercase Unicode accent normalization, sublinear TF, smoothed IDF,
and L2 normalization. Missing text is empty. LR uses $C=1$, liblinear, 500 iterations, and
tolerance $10^{-4}$. Reason labels and reason embeddings are excluded from outcome predictors.

Separate-field models use criteria/disease/drug vocabulary limits 20,000/5,000/5,000 and block
scaling $1/\sqrt{3}$. Character models use word-boundary 3--5 grams, minimum document frequency 3,
and at most 30,000 features per phase. The historical blend weights concatenated word,
separate-field word, and character probabilities 0.25/0.25/0.50. Text plus metadata and pathway
features did not improve validation loss: the $C=1$ metadata model scored 0.59460, compared with
0.58377 for text alone. Exact drug matching and sparse pathway coverage can limit added signal;
these are plausible explanations rather than isolated causal effects.

\subsection{Phase sharing and later cohorts}
The new text LR component shares coefficients across phases, adds phase-specific interactions
scaled by 0.35 and phase intercepts, and uses training recency weights with a six-year half-life.
Start year changes training weights, not inference features. This permits common text effects
while retaining phase differences. Its final probability receives weight 0.5 alongside the
historical blend.

A historical screen trains through 2009 and evaluates 2010--2013; equal-phase loss improves
from 0.58425 to 0.57288 for the selected blend. Five trial-grouped training folds select the
mixture; OOF loss improves from 0.59649 to 0.59155. The supplied validation check then improves
from 0.58090 to 0.56839. These comparisons motivated the public submission, which scores
0.59011. Outcome training does not add the supplied validation labels.

\subsection{Overall performance and generalization}
Across all 1,013 validation trial-phase rows, pooled F1/ROC-AUC/AP improve from
0.7354/0.7569/0.7982 to 0.7387/0.7737/0.8144. These pooled metrics differ from the
equal-phase mean log loss used for selection. Lower loss need not improve every phase's F1.

Public log loss is 0.59011 versus validation 0.56839, a gap of 0.02172.
Later trial cohorts, repeated validation selection, and public-subset sampling variability
may contribute; this comparison does not isolate any one cause. Private performance is unknown.

\clearpage
\twocolumn[{
\section{Failure Errors by Phase and Reason}
\begin{center}\small
\begin{tabular}{llrrll}
\toprule
Phase & Reason group & $n$ & Obs./Pred. & Baseline errors & Current blend errors \\
\midrule
I & Safety & 10 & 1/9 & 4/10 (40.0\%) & 5/10 (50.0\%) \\
I & Funding & 4 & 1/3 & 2/4 (50.0\%) & 2/4 (50.0\%) \\
I & Efficacy & 5 & 0/5 & 5/5 (100.0\%) & 5/5 (100.0\%) \\
I & Recruitment & 22 & 2/20 & 10/22 (45.5\%) & 7/22 (31.8\%) \\
I & Other & 2 & 0/2 & 1/2 (50.0\%) & 1/2 (50.0\%) \\
I & No recorded reason & 28 & -- & 17/28 (60.7\%) & 14/28 (50.0\%) \\
\midrule
II & Safety & 12 & 1/11 & 5/12 (41.7\%) & 5/12 (41.7\%) \\
II & Funding & 4 & 1/3 & 0/4 (0.0\%) & 0/4 (0.0\%) \\
II & Efficacy & 18 & 0/18 & 4/18 (22.2\%) & 5/18 (27.8\%) \\
II & Recruitment & 91 & 12/79 & 17/91 (18.7\%) & 17/91 (18.7\%) \\
II & Other & 15 & 3/12 & 5/15 (33.3\%) & 4/15 (26.7\%) \\
II & No recorded reason & 144 & -- & 48/144 (33.3\%) & 48/144 (33.3\%) \\
\midrule
III & Safety & 10 & 1/9 & 7/10 (70.0\%) & 7/10 (70.0\%) \\
III & Funding & 0 & 0/0 & -- & -- \\
III & Efficacy & 10 & 0/10 & 10/10 (100.0\%) & 10/10 (100.0\%) \\
III & Recruitment & 23 & 5/18 & 14/23 (60.9\%) & 14/23 (60.9\%) \\
III & Other & 6 & 2/4 & 5/6 (83.3\%) & 5/6 (83.3\%) \\
III & No recorded reason & 82 & -- & 52/82 (63.4\%) & 50/82 (61.0\%) \\
\bottomrule
\end{tabular}

\captionof{table}{True outcome failures predicted successful at threshold 0.5: $k/n$ and percentage.
Both models use the same historical Part 1 reason lookup. ``Obs./Pred.'' counts observed versus
estimated reason memberships; groups overlap. Unmatched trials form the no-recorded-reason
group. A dash indicates no eligible trials or no applicable reason source.}
\label{tab:errors}
\end{center}
}]
\subsection{Definition and provenance}
For each phase/reason group $g$ containing only true failures,
\begin{equation}
 \mathrm{FSR}_{g}=\frac{\sum_{i\in g}\mathbb{I}[p_i\geq0.5]}{|g|}.
\end{equation}
A trial counts once within each phase/group, but may belong to several groups.
An absent Part 1 record is not a negative reason label.

Freezing the historical K-Means/LR reason lookup isolates changes in outcome predictions.
This is not an analysis of the newer 235-label ridge reason model. We distinguish observed,
estimated, and training-resubstitution reasons.

The 452 failed trial-phase rows represent 398 unique trials. Of 198 matched rows, 27 have observed reasons and 171 have predicted reasons; 254 rows have no Part 1 match. Among predicted rows, 131 are training resubstitution estimates. These are not out-of-fold or confirmed labels. Matched empty-reason groups have zero support after fallback and are omitted.


\subsection{Interpretation}
Phase I false-success counts change from 34/64 to 29/64, while true-success false negatives change from 9 to 13.

Phase II false-success counts change from 73/265 to 73/265, while true-success false negatives change from 96 to 94.

Phase III false-success counts change from 80/123 to 78/123, while true-success false negatives change from 21 to 20.

Phase I Recruitment errors fall from 10/22 to 7/22, while Safety errors rise from 4/10 to 5/10. Phase II Other improves from 5/15 to 4/15, but Efficacy worsens from 4/18 to 5/18. All five Phase I and all ten Phase III Efficacy failures remain false successes. These groups include estimated reasons.


Small groups make one-trial changes look large; Phase III has no Funding failures.
Estimated reasons add uncertainty beyond sampling variability. These descriptive associations
do not establish causes of errors. Stopping reasons never enter the outcome model as features.

\clearpage
\twocolumn[{
\section{Discussion and Reproducibility}
\begin{center}\small
\begin{tabular}{p{.25\textwidth}p{.42\textwidth}p{.24\textwidth}}
\toprule
Additional experiment family & Main evidence & Decision\\
\midrule
Graph~\cite{graph} / soft or selective pseudo labels & Graph validation F1 0.8932; public F1 0.84081 versus incumbent 0.85405. & Negative generalization example.\\
Primary/secondary reasons and whitening & Primary/secondary OOF 0.8960, validation 0.8590; shrinkage whitening did not pass checks. & Retain summary.\\
Label shift~\cite{saerens,bbse}; residual Other/text & Prior correction gains one late trial; residual Other worsens late-year F1. & No promotion.\\
Local ensemble, Other kernel, F1 decoder & Fixed-ridge OOF 0.9026; sequential procedure 0.9002. & Retain public 0.88568.\\
Ridge grid and group weighting & Grid-selected OOF 0.8983; weighting changes no final OOF decisions. & Retain public 0.88568.\\
\bottomrule
\end{tabular}
\captionof{table}{Later unsuccessful hypotheses. Earlier 157-label OOF protocols differ from
the newer 235-label near-group protocol and cannot be ranked directly. Fixed-recipe studies
followed recipe selection. Sources and recorded evidence are included in the code ZIP.}
\label{tab:negative}
\end{center}
}]
\subsection{Why more features need not help Part 1}
The embedding-only choice matches the target: stopping-reason text can express recruitment,
funding, or administrative problems that a drug's biology does not describe.
The pathway decline in Section~\ref{sec:p1features} supports excluding that tested block.
Possible explanations are incomplete exact-name drug matching (about 56\% coverage in Part 1~\cite{handout}),
coarse or redundant pathways, and extra coefficients estimated from only 157 gold labels;
pseudo targets add uncertain supervision, not 1,014 independent gold labels.
Missing drug matches also do not imply biological absence. These are plausible mechanisms,
not separately established causes of the observed decline.

Criteria, disease, drug, ICD, target/action, and molecular representations can increase variance
or dilute reason similarity. No exhaustive all-feature or fingerprint ablation was run;
their usefulness is unresolved.
For the chosen kernel, global scaling preserves Euclidean geometry while column-wise scaling
alters it. The public improvement also changes model family and gold count (157 to 235),
so it cannot isolate geometry. Local models avoid inheriting uncertain cluster targets;
rare-label suppression is a limitation of the tested K-Means/GMM fits, not all semi-supervised methods.

Follow-ups did not improve the pipeline. The sample-F1 decoder assumes independent Bernoulli
marginals, potentially violated by correlated reasons. A 16-setting ridge grid and inverse-group
weighting also failed: 82/235 labels occupy 13 repeated groups, mostly Recruitment,
but downweighting did not improve held-group decisions.

\subsection{Selection and temporal limitations}
Validation labels and public feedback were inspected repeatedly. Later nested resampling
controls leakage within each fit but cannot erase that search history. The final Part 1 fit
uses all 157+78 labels after recipe selection. Historical validation scores come from separate
157-label fits; they are not independent scores of the final 235-label artifact.

The late-20 and late-56 diagnostics have different training populations and reason mixtures:
late train contains one Recruitment-positive trial, versus 21 in late validation. Neither
small subset determines future performance. Temporal fits retain earlier duplicates available at that date; saved diagnostics
isolate novel groups. Test data never enter fitting or selection. Private performance is unknown.

\subsection{Reproduction and preserved evidence}
The report uses CVPR style~\cite{template}. \path{models/selection.json} fixes final hashes.
Run \texttt{python reproduce.py} from the project root to verify saved predictions and IDs.
The ZIP includes data/audits, models/evidence, source/tests/notebooks, and report/LaTeX in
\texttt{data/}, \texttt{models/}, \texttt{tools/}, and \texttt{report/}, respectively.
Run \path{report/scripts/report_analysis.py} then \path{report/scripts/generate_tables.py}
with Python to regenerate metrics, figures, and tables. Five notebooks in \path{tools/notebooks/}
replay saved evidence with training disabled. Historical runners/manifests are included;
full backups and the frozen November 1 delivery remain local.

\paragraph{Team contributions and assistance.}
Pei-Chen Chung (pc786) and Yu-Jye Hu (yh2444) contributed equally to data preprocessing,
exploratory data analysis, feature engineering, model development, performance evaluation,
and preparation of the final report. Both members shared responsibilities, jointly discussed
key modeling decisions, and contributed equally to the final deliverables. OpenAI Codex
assisted with implementation, verification, organization, and drafting. Both authors are
responsible for the submitted work.

\clearpage
{\small
\begin{thebibliography}{99}
\raggedright
\bibitem{handout} CS5785 / ORIE5750 / ECE5414 teaching staff.
\emph{Midterm Project 2026}. Cornell University, Fall 2026. Course assignment and supplied datasets.
\bibitem{sklearn} F. Pedregosa et al. Scikit-learn: Machine Learning in Python.
\emph{Journal of Machine Learning Research}, 12:2825--2830, 2011.
\url{https://jmlr.org/papers/v12/pedregosa11a.html}.
\bibitem{tfidf} Scikit-learn developers. \emph{Text feature extraction}.
\url{https://scikit-learn.org/stable/modules/feature_extraction.html}.
\bibitem{lectures} CS5785 / ORIE5750 / ECE5414 teaching staff.
\emph{Classification Evaluation, Regularization, Model Selection, K-Means, GMM, and EM}.
Supplied course lectures, Fall 2026.
\bibitem{graph} Scikit-learn developers. \emph{LabelSpreading}.
\url{https://scikit-learn.org/stable/modules/generated/sklearn.semi_supervised.LabelSpreading.html}.
\bibitem{svm} Scikit-learn developers. \emph{Support vector machines}.
\url{https://scikit-learn.org/stable/modules/svm.html}.
\bibitem{normalizer} Scikit-learn developers. \emph{Normalizer}.
\url{https://scikit-learn.org/stable/modules/generated/sklearn.preprocessing.Normalizer.html}.
\bibitem{template} CVPR organizers. \emph{CVPR 2026 Author Kit}, commit
\texttt{12909ae437f6}.
The unmodified style file is retained with this report.
\url{https://github.com/cvpr-org/author-kit}.
\bibitem{saerens} M. Saerens, P. Latinne, and C. Decaestecker.
Adjusting the Outputs of a Classifier to New a Priori Probabilities: A Simple Procedure.
\emph{Neural Computation}, 14(1):21--41, 2002.
\url{https://doi.org/10.1162/089976602753284446}.
\bibitem{bbse} Z. Lipton, Y.-X. Wang, and A. Smola.
Detecting and Correcting for Label Shift with Black Box Predictors.
\emph{Proceedings of ICML}, PMLR 80:3122--3130, 2018.
\url{https://proceedings.mlr.press/v80/lipton18a.html}.
\end{thebibliography}

\medskip
\noindent\textbf{Code package.}
The complete code ZIP (\texttt{midterm\_code.zip}) is available at
\url{https://drive.google.com/file/d/1BBR6opsUrpWm2WhqhjemkzKJEbT6cnHd/view?usp=sharing}.
}

\end{document}
